\begin{figure}[!htbp]
\centering
\resizebox{.76\linewidth}{!}{%
\begin{tikzpicture}[scale=3.0,>=Latex]
  \HexCoords
  \DrawHexagon
  \DrawRays

  \foreach \i in {0,...,5}{
    \coordinate (Pr\i) at ($(O)!.16!(V\i)$);
    \fill[DemandRed] (Pr\i) circle (.018);
  }
  \draw[DemandRed!65,line width=.65pt]
    (Pr0)--(Pr1)--(Pr2)--(Pr3)--(Pr4)--(Pr5)--cycle;
  \draw[CenterBlue,dashed,fill=CenterBlue!7,line width=.75pt]
    (O) circle ({.16*sqrt(3)/2});

  \coordinate (Qm) at (-.18,-.50);
  \coordinate (Qz) at (-.02,-.62);
  \coordinate (Qp) at (.17,-.52);
  \fill[GapPurple] (Qm)--++(.025,.043)--++(.025,-.043)--cycle;
  \fill[GapPurple] (Qz)--++(.025,.043)--++(.025,-.043)--cycle;
  \fill[GapPurple] (Qp)--++(.025,.043)--++(.025,-.043)--cycle;
  \node[figlabel,left=2pt] at (Qm) {$Q_-$};
  \node[figlabel,below=2pt] at (Qz) {$Q_0$};
  \node[figlabel,right=2pt] at (Qp) {$Q_+$};
  \node[figlabel,right] at (Pr0) {$P_0^{\rm rad}$};
  \node[figlabel,above right] at (Pr1) {$P_1^{\rm rad}$};
  \node[figlabel,below left] at (Pr4) {$P_4^{\rm rad}$};
  \fill (O) circle (.012) node[above left=1pt,figlabel] {$O$};

  \draw[centerrole]
    (-.49,-.73)--(.50,-.72)--(.02,.24)--cycle;
  \node[figlabel,text=CenterBlue] at (.43,-.63) {putative $U_C$};

  \node[figlabel,atlasbox,align=left,anchor=west] at (1.20,.55)
    {six symmetric points: $P_i^{\rm rad}=(1-c_*)V_i$\\
     three asymmetric points: $Q_-,Q_0,Q_+$\\
     $\mathbb B(0,r_*)\subset\operatorname{conv}\{P_i^{\rm rad}\}$};
  \draw[flow] (1.15,.30)--(.34,.05);

  \node[figlabel,align=center] at (0,-1.22)
    {$K_{\rm wit}=\mathbb B(0,r_*)\cup\{Q_-,Q_0,Q_+\}
      \subset U_C$ under a hypothetical cover};
\end{tikzpicture}%
}
\caption{The zero-gap nine-point witness.  The six symmetric radial points
force a centered disk into the convex C triangle, while the three asymmetric
frontier points supply the missing directional support.  Marker shapes
separate the two constructions: circles are radial witnesses and triangles
are strict-supercritical $AB$-frontier witnesses.  The displayed placement is
illustrative; the point definitions are the exact ones in the preceding
lemmas.}
\label{fig:zero-gap-nine-point-witness}
\end{figure}
